\documentclass{beamer}
\usepackage{amsmath, amsfonts, amssymb, array, graphicx, mathtools, tabularx}
\usepackage[english]{babel}
\usepackage[utf8]{inputenc}

\usetheme{Madrid}

\title{VNS}
\author{Hugo Diaz Delgado \and Alessandro Lemme}
\institute{Università degli studi di Brescia}
\date{\today}

\begin{document}

\begin{frame}
\titlepage
\end{frame}

\begin{frame}{Outline}
\tableofcontents
\end{frame}

\section{Problem definition}

\subsection{A little notation}
\begin{frame}{What are we doing}
\quad Let's define a minimization problem:
\begin{gather*}
    \min f(x) \quad \text{subject to} \quad x \in X
\end{gather*}
\quad Given the set of feasible solutions $X$, we want to find $x^{*} \in X$ such that 
\begin{gather*}
    f(x^{*}) \leq f(x) \ \forall x \in X
\end{gather*} 
\quad If such result exists, an \textbf{exact algorithm} also proves its optimality (or shows $X = \emptyset$). If dealing with a continuous function, we have to admit a certain degree of tolerance such that
\begin{gather*}
    \frac{f(x) - f(x^{*})}{f(x^{*})} < \varepsilon \ \forall x \in X
\end{gather*} 
with $\varepsilon$ small and positive.
\end{frame}

\begin{frame}{Why heuristics?}
\quad Many instances in Operational Research are \textbf{NP-hard} (or have high polynomial complexity), so finding $x^{*}$ with exact methods in a reasonable time is impossible. \\
\vspace{0.2cm}
\quad One is often forced to resort to \textbf{heuristics}, which find a solution $x_h$ "good enough" quickly, trading efficacy for efficiency (without proof of optimality). \\
\vspace{0.2cm}
\quad Some heuristics have a worst-case guarantee:
\begin{gather*}
    \frac{f(x_{h}) - f(x)}{f(x_{h})} \leq \varepsilon \ \forall x \in X
\end{gather*}
although this theoretical $\varepsilon$ is usually much larger than the error observed in practice.
\end{frame}

\begin{frame}{Local optima}
\quad Given a neighborhood $N(x)$ of $x$, we define a local optimum $x_{L}$ as
\begin{gather*}
    x_{L} \in X : f(x_{L}) \leq f(x) \ \forall x \in N(x_{L}) \cap X
\end{gather*} 
however $f(x^{*})$ could differ substantially from the average value of a local minimum. There are many ways to get out of a local optimum and its valley. \\
\vspace{0.2cm}
\textbf{Metaheuristics} are a general framework for combinatorial and global optimization problems (e.g., Simulated Annealing). \\
\vspace{0.2cm}
\textbf{Variable Neighborhood Search (VNS)} is a metaheuristic that exploits systematically the idea of neighborhood change, both to descend to local minima and to escape from the valleys containing them. Moreover, it requires \textbf{few or no parameters} and gives insight into its performance.
\end{frame}

\begin{frame}{Observations}
VNS systematically exploits three facts:
\vspace{0.2cm}
\begin{itemize}
\item A local minimum in a neighborhood is not necessarily so for another.
\item A global minimum is a local minimum in every possible neighborhood.
\item For many problems, local minima with respect to one or several neighborhoods are relatively close to each other.
\end{itemize}

\vspace{0.3cm}
Corollary of the empirical observation 3: a local optimum gives information about the global one (e.g., sharing variable values), so an organized study of its neighborhood is useful.
\end{frame}

\begin{frame}{Tutorial structure}
\begin{itemize}
\item \textbf{Preliminaries:} Gathering and evaluating information about the problem.
\item \textbf{VND (Variable Neighborhood Descent):} Deterministic descent.
\item \textbf{RVNS (Reduced VNS):} Stochastic neighborhood exploration.
\item \textbf{Basic and General VNS:} Merging VND and RVNS.
\item \textbf{SVNS (Skewed VNS):} Getting out of very large valleys.
\item \textbf{VNDS (Variable Neighborhood Decomposition Search):} Two-level scheme for very large instances.
\item \textbf{Performance analysis:} Distance-to-target diagrams and valley profiles.
\end{itemize}
\end{frame}




\section{Preliminaries: Documentation}

\begin{frame}{Literature review}
\quad Before designing a VNS, we must collect existing papers about the problem (which may have different names in different fields).

\vspace{0.2cm}
We study the literature to answer four main questions:
\begin{enumerate}
\item \textbf{Difficulty:} Is the problem NP-hard? If it is in P, is the exact algorithm fast enough for real instances?
\item \textbf{Exact algorithms:} Are there test instances available? What is the largest size solved exactly?
\item \textbf{Previous heuristics:} Which heuristics already exist, and how good are they (size solved, error, and time)?
\item \textbf{Algorithmic steps:} What moves and neighborhoods are used? Is there existing code we can reuse?
\end{enumerate}
\end{frame}

\begin{frame}{Why do we need this information?}
\begin{itemize}
\item \textbf{Question 1 (Justification):} Tells us if we really need a VNS heuristic.
\vspace{0.2cm}
\item \textbf{Questions 2 \& 3 (Benchmarking):} Give us a reference to test our VNS:
    \begin{itemize}
    \item It must find the optimal solution for almost all instances solved by exact methods.
    \item To be useful, it must match previous heuristics and perform \textbf{much better} on several instances.
    \end{itemize}
\vspace{0.2cm}
\item \textbf{Question 4 (Ingredients \& Reuse):} Helps us choose neighborhoods for VND, find unexplored ideas, and reuse standard code.
\end{itemize}
\end{frame}

%  a partir de aqui todo esto es del punto 3


\section{Variable Neighborhood Descent (VND)}

\subsection{Local Search Basics}
\begin{frame}{Descent heuristics}
\quad Before talking about VND, let's define two basic local search methods starting from an initial solution $(x)$ within a neighborhood $N(x)$:

\vspace{0.3cm}
\begin{itemize}
\item \textbf{Steepest descent (Best improvement):} Checks the entire neighborhood $N(x)$ and moves to the best neighbor $x' = \arg\min_{y \in N(x)} f(y)$ if $f(x') < f(x)$.
\vspace{0.2cm}
\item \textbf{First descent (First improvement):} Checks neighbors one by one and moves to the first neighbor $x'$ that improves $f(x)$. It is faster when $N(x)$ is big.
\end{itemize}
\end{frame}

\subsection{The VND Method}
\begin{frame}{VND: Idea and Objective}
\quad Basic descent methods use only one neighborhood, so they stop at the first local minimum. \textbf{Variable Neighborhood Descent (VND)} solves this by combining several neighborhoods in a deterministic way.

\vspace{0.2cm}
\quad It is based on a simple idea: a local minimum for one neighborhood $N_1(x)$ is not necessarily a local minimum for another neighborhood $N_2(x)$.

\vspace{0.3cm}
\textbf{Main objective:}
\begin{itemize}
\item Escape local optima by changing to a new neighborhood $N_\ell$ when the search gets stuck.
\item Find a better solution that is a local optimum for all neighborhoods ($N_1, \dots, N_{\ell_{\max}}$).
\end{itemize}
\end{frame}

\begin{frame}{How Basic VND works}
\textbf{Initialization:} Choose a set of neighborhoods $N_\ell$ ($\ell = 1, \dots, \ell_{\max}$) and find an initial solution $x$.

\vspace{0.2cm}
\textbf{Main loop} (repeat until no improvement is found):
\begin{enumerate}
\item Set $\ell \leftarrow 1$.
\item Repeat until $\ell = \ell_{\max}$:
    \begin{itemize}
    \item \textbf{Explore:} Find the best neighbor $x' \in N_\ell(x)$.
    \item \textbf{Move or change:} If $f(x') < f(x)$, set $x \leftarrow x'$ and go back to $\ell \leftarrow 1$; otherwise, try the next neighborhood $\ell \leftarrow \ell + 1$.
    \end{itemize}
\end{enumerate}
\end{frame}

\subsection{Design Considerations}
\begin{frame}{Practical considerations}
\quad To make VND work well, we need to think about three things:

\vspace{0.3cm}
\begin{itemize}
\item \textbf{Order \& Speed:} Order neighborhoods from smallest (fastest) to largest (slowest), going back to $\ell = 1$ after every improvement to save time.
\vspace{0.2cm}
\item \textbf{Exploration:} Simple moves might not be enough to escape narrow valleys, so we may need bigger moves.
\vspace{0.2cm}
\item \textbf{Precision:} We can use all the potential of VND if we use it in or with some other method like VNS. This mix could give us better and faster solutions.
\end{itemize}
\end{frame}


%   a partir de aqui estamos en la seccion 4

\section{Reduced Variable Neighborhood Search (RVNS)}

\subsection{The RVNS Method}
\begin{frame}{RVNS: Idea and Key Questions}
\quad Now we will talk about, \textbf{Reduced VNS (RVNS)} which is a pure \textbf{probabilistic} method. Instead of doing a full local search, it picks random points from nested neighborhoods to escape local minima.

\vspace{0.3cm}
To leave a valley without prior knowledge, RVNS answers three questions:
\begin{itemize}
\item \textbf{Which direction to go?} Choose a direction at random, so any point can be reached.
\vspace{0.15cm}
\item \textbf{How far?} Start close to $(x)$, due to the fact that good local optima tend to be near each other.
\vspace{0.15cm}
\item \textbf{What if moves fail?} Go further by switching to a larger neighborhood $N_{k+1}(x)$.
\end{itemize}
\end{frame}

\begin{frame}{How RVNS works}
\textbf{Initialization:} Choose nested neighborhood structures $N_k$ ($k = 1, \dots, k_{\max}$), an initial solution $x$, and a stopping condition (e.g., max time or iterations).

\vspace{0.2cm}
\textbf{Main loop} (repeat until the stopping condition is met):
\begin{enumerate}
\item Set $k \leftarrow 1$.
\item Repeat until $k = k_{\max}$:
    \begin{itemize}
    \item \textbf{Shaking:} Pick a random point $x' \in N_k(x)$.
    \item \textbf{Move or not:} If $f(x') < f(x)$, move there ($x \leftarrow x'$) and go back to $N_1$ ($k \leftarrow 1$); otherwise, try a larger neighborhood ($k \leftarrow k + 1$).
    \end{itemize}
\end{enumerate}
\end{frame}

\begin{frame}{Why use RVNS?}
    \begin{itemize}
        \item \textbf{Speed:} By omitting the deterministic local search phase, each iteration evaluates only a single random point, making it extremely fast.
        \vspace{0.3cm}
        \item \textbf{Focused exploration:} Nested neighborhoods ($\mathcal{N}_1(x) \subset \mathcal{N}_2(x) \subset \dots \subset \mathcal{N}_{k_{\max}}(x)$) analyze the closest neighborhoods of $(x)$ more intensively before jumping to distant regions.
        \vspace{0.3cm}
        \item \textbf{Role in VNS:} Ideal when local descent is computationally too expensive, or as a fast stochastic initialization step to improve the starting solution in General VNS.
    \end{itemize}
\end{frame}



%a partir de aqui esto es la seccion 5 

\section{Basic and General VNS}

\subsection{Combining VND and RVNS}
\begin{frame}{From VND and RVNS to VNS}
\quad In the previous sections, we saw two complementary tools:
\begin{itemize}
\item \textbf{VND:} Uses deterministic local search to find good local optima, but can still get trapped in deep valleys.
\item \textbf{RVNS:} Uses random shaking to escape valleys fast, but does not refine solutions with local descent.
\end{itemize}

\vspace{0.3cm}
\quad \textbf{Variable Neighborhood Search (VNS)} combines both ideas:
\begin{itemize}
\item \textbf{Basic VNS:} Combines random shaking with a simple local search.
\item \textbf{General VNS (GVNS):} Improves the initial solution with \textbf{RVNS} and uses \textbf{VND} as the local search method.
\end{itemize}
\end{frame}

\subsection{Basic VNS}
\begin{frame}{How Basic VNS works}
\textbf{Initialization:} Choose neighborhood structures $\mathcal{N}_k$ ($k = 1, \dots, k_{\max}$), an initial solution $x$, and a stopping condition.

\vspace{0.2cm}
\textbf{Main loop} (repeat until the stopping condition is met):
\begin{enumerate}
\item Set $k \leftarrow 1$.
\item Repeat until $k = k_{\max}$:
    \begin{itemize}
    \item \textbf{Shaking:} Pick a random point $x' \in \mathcal{N}_k(x)$.
    \item \textbf{Local search:} Apply a local search starting from $x'$ to get a local optimum $x''$.
    \item \textbf{Move or not:} If $f(x'') < f(x)$, move there ($x \leftarrow x''$) and reset $k \leftarrow 1$; otherwise, set $k \leftarrow k + 1$.
    \end{itemize}
\end{enumerate}
\end{frame}

\begin{frame}{The three possible outcomes}
\quad After applying local search to $x'$, we obtain a local minimum $x''$. There are three possible cases:

\vspace{0.3cm}
\begin{enumerate}
\item \textbf{$x'' = x$:} We fell back to the bottom of the same valley. We continue with the next neighborhood ($k \leftarrow k + 1$).
\vspace{0.2cm}
\item \textbf{$x'' \neq x$ and $f(x'') \geq f(x)$:} We found a different local minimum, but it is not better than our current best solution. We continue with $k \leftarrow k + 1$.
\vspace{0.2cm}
\item \textbf{$x'' \neq x$ and $f(x'') < f(x)$:} We found a better local minimum. We update $x \leftarrow x''$ and restart from the first neighborhood ($k \leftarrow 1$).
\end{enumerate}
\end{frame}

\subsection{General VNS and Neighborhood Design}
\begin{frame}{General VNS (GVNS)}
\quad \textbf{General VNS} upgrades Basic VNS by using variable neighborhoods in both phases:

\vspace{0.2cm}
\textbf{Initialization:} Select shaking neighborhoods $\mathcal{N}_k$ ($k = 1, \dots, k_{\max}$) and descent neighborhoods $N_\ell$ ($\ell = 1, \dots, \ell_{\max}$). Find an initial solution $x$ and improve it with \textbf{RVNS}.

\vspace{0.2cm}
\textbf{Main loop} (for $k = 1, \dots, k_{\max}$):
\begin{enumerate}
\item \textbf{Shaking:} Generate a random point $x' \in \mathcal{N}_k(x)$.
\item \textbf{Local search by VND:} Apply VND starting from $x'$ using $N_1, \dots, N_{\ell_{\max}}$ to obtain $x''$.
\item \textbf{Move or not:} If $f(x'') < f(x)$, set $x \leftarrow x''$ and $k \leftarrow 1$; otherwise, $k \leftarrow k + 1$.
\end{enumerate}
\end{frame}

\begin{frame}{Neighborhood properties}
\quad To avoid getting trapped and ensure we can reach the global optimum, neighborhoods should satisfy:

\vspace{0.2cm}
\begin{itemize}
\item \textbf{Cover the feasible space:} From any solution $x \in X$, the union of neighborhoods should cover $X$:
\begin{gather*}
    X \subseteq \mathcal{N}_1(x) \cup \mathcal{N}_2(x) \cup \dots \cup \mathcal{N}_{k_{\max}}(x) \ \forall x \in X
\end{gather*}
\item \textbf{Nested structures:} Usually built by repeating a basic move $k$ times, so sizes increase:
\begin{gather*}
    \mathcal{N}_1(x) \subset \mathcal{N}_2(x) \subset \dots \subset \mathcal{N}_{k_{\max}}(x), \quad X \subset \mathcal{N}_{k_{\max}}(x)
\end{gather*}
This explores close regions more thoroughly before jumping far.
\end{itemize}
\end{frame}

% a partir de aqui va a ser la seccion 7

\section{Skewed Variable Neighborhood Search (SVNS)}

\subsection{Motivation and Idea}
\begin{frame}{Skewed VNS?}
\quad Standard VNS works best when good local optima are close to each other. However, some problems have \textbf{several deep valleys located far apart}.

\vspace{0.2cm}
In those cases, standard VNS has two issues:
\begin{itemize}
\item Using very large neighborhoods makes it behave like a random \textit{multistart}.
\item If the current solution $x$ is not in the deepest valley, staying strictly near it prevents finding better distant regions.
\end{itemize}

\vspace{0.2cm}
\textbf{Skewed VNS (SVNS)} solves this by allowing the search to move to a new valley even if its local minimum is slightly worse than $x$, as long as it is \textbf{far away}.
\end{frame}

\subsection{The Decision Criterion and Algorithm}
\begin{frame}{When is it worth moving?}
\quad To decide if moving to a new local optimum $x''$ is worth it, SVNS relaxes the acceptance rule by adding a distance bonus:
\begin{gather*}
    f(x'') - \alpha \rho(x, x'') < f(x)
\end{gather*}
where:
\begin{itemize}
\item $\rho(x, x'')$ is the \textbf{distance} between the current solution $x$ and $x''$.
\item $\alpha > 0$ is a parameter that controls the compensation between solution quality and distance.
\end{itemize}

\vspace{0.2cm}
$\Rightarrow$ If $x''$ is far enough from $x$, we accept it to explore that new valley, even if $f(x'') \geq f(x)$.
\end{frame}

\begin{frame}{How SVNS works}
\textbf{Initialization:} Choose neighborhoods $\mathcal{N}_k$ ($k = 1, \dots, k_{\max}$), an initial solution $x$, set $x_{opt} \leftarrow x$, $f_{opt} \leftarrow f(x)$, and choose $\alpha$.

\vspace{0.5cm}
\textbf{Main loop} (for $k = 1, \dots, k_{\max}$):
\vspace{0.1cm}
\begin{enumerate}
\item \textbf{Shaking:} Generate a random point $x' \in \mathcal{N}_k(x)$.
\item \textbf{Local search:} Apply local search from $x'$ to obtain $x''$.
\item \textbf{Save best:} If $f(x'') < f_{opt}$, update $x_{opt} \leftarrow x''$ and $f_{opt} \leftarrow f(x'')$.
\item \textbf{Move or not:} If $f(x'') - \alpha \rho(x, x'') < f(x)$, move to $x \leftarrow x''$ and reset $k \leftarrow 1$; otherwise, set $k \leftarrow k + 1$.
\end{enumerate}

\vspace{0.15cm}
\end{frame}


% a partir de aqui todo es seccion 7


\section{Variable Neighborhood Decomposition Search (VNDS)}

\subsection{Introduction and Idea}
\begin{frame}{What is VNDS?}
\quad When problem instances are very large, exploring the entire solution space $S$ at every step becomes too slow. \textbf{Variable Neighborhood Decomposition Search (VNDS)} extends Basic VNS into a \textbf{two-level scheme} based on problem decomposition.

\vspace{0.3cm}
\textbf{Main objectives:}
\begin{itemize}
\item Break large, complex problems into smaller, manageable subproblems.
\item Apply local search only to a small part of the solution while keeping the rest fixed, saving computing time.
\end{itemize}

\vspace{0.3cm}
\end{frame}

\subsection{Comparison and Algorithm}
\begin{frame}{VNDS vs Basic VNS: How it works}
\textbf{Key difference (Step 2b):} Basic VNS runs local search in the \textbf{whole space} $S$, while VNDS solves a subproblem only inside a \textbf{subspace} $V_k \subseteq \mathcal{N}_k(x)$.

\vspace{0.25cm}
\textbf{Main loop} (starting from an initial solution $x$, for $k = 1, \dots, k_{\max}$):
\begin{enumerate}
\item \textbf{Shaking:} Pick $x' \in \mathcal{N}_k(x)$. Let $y = x' \setminus x$ be the $k$ new attributes in $x'$.
\item \textbf{Subspace local search:} Optimize only over $y$ to get $y'$, and rebuild the full solution:
\begin{gather*}
    x'' = (x' \setminus y) \cup y'
\end{gather*}
\item \textbf{Move or not:} If $f(x'') < f(x)$, set $x \leftarrow x''$ and $k \leftarrow 1$; otherwise, set $k \leftarrow k + 1$.
\end{enumerate}
\end{frame}


%esto es punto 8

\section{Performance Analysis}

\begin{frame}{Diagnostic tools}
\quad To understand how VNS behaves and tune its design, we use two visual diagnostic tools:

\vspace{0.25cm}
\begin{itemize}
\item \textbf{Distance-to-target diagrams:} Plot the distance from the current solution to the best-known optimum at each iteration. They show whether the search is moving closer to the target or getting trapped.
\vspace{0.25cm}
\item \textbf{Mountain and valley profiles:} Built by moving step-by-step along a path between two local optima and plotting $f(x)$. They reveal the \textbf{width and height of the barriers} separating valleys:
    \begin{itemize}
    \item If barriers are low and optima are clustered, standard VNS works well.
    \item If valleys are very wide and separated by high mountains, we need larger neighborhoods or \textbf{SVNS}.
    \end{itemize}
\end{itemize}
\end{frame}


%  indice de cosas

\section{Conclusions}

\begin{frame}{Conclusions: Why VNS works}
\quad VNS satisfies the eight desirable properties of a good metaheuristic:

\vspace{0.25cm}
\begin{itemize}
\item \textbf{Simplicity \& Coherence:} Based on a clear, general principle (systematic neighborhood change) from which all steps naturally follow.
\vspace{0.15cm}
\item \textbf{Precision:} Formulated in exact mathematical terms, without relying on physical or biological metaphors.
\vspace{0.15cm}
\item \textbf{Efficiency \& Effectiveness:} Finds optimal or near-optimal solutions in moderate computing time.
\vspace{0.15cm}
\item \textbf{Robustness:} Performs consistently across many different problem instances, without needing heavy fine-tuning.
\vspace{0.15cm}
\item \textbf{User-friendliness \& Innovation:} Easy to understand and use (requires very few or no parameters), enabling new applications.
\end{itemize}
\end{frame}


\end{document}